\documentclass[11pt]{article}
\usepackage[margin=1in]{geometry}
\usepackage{setspace}
\usepackage{hyperref}
\hypersetup{hidelinks}
\setstretch{1.08}

\title{Response to the Reviewer\\[0.4em]
\large Manuscript JCP26-AR-03161}
\author{Vincenzo Barone}
\date{}

\begin{document}
\maketitle

Dear Editor and Reviewer,

Thank you for the careful and positive evaluation of our manuscript,
``Algorithmic Construction and Exploitation of Topology-Constrained Refinement
Models for Semiexperimental Equilibrium Structures.'' We respond below to the
reviewer's comments in full. Textual changes are shown in red in the marked
manuscript. Changes to figure artwork are not colored red, but are described
explicitly in the corresponding responses.

\section*{Reviewer \#1}

\subsection*{General assessment}

\paragraph{Reviewer comment.}
\emph{The manuscript ``Algorithmic Construction and Exploitation of
Topology-Constrained Refinement Models for Semiexperimental Equilibrium
Structures'' introduces a significant methodological advancement in the field
of semiexperimental equilibrium structure refinement. The determination of
$r_e^{\mathrm{SE}}$ structures has traditionally been an ``art'' reserved for
experts due to the complexities of coordinate selection, symmetry handling,
and rank deficiency. MORPHEUS addresses these issues by providing a
reproducible, graph-driven framework that automates model construction. The
validation on systems ranging from glycine to a 49-atom testosterone molecule
demonstrates its robustness. The main strengths are the reproducibility,
handling of incomplete data, validation and clarity of the results and work.
This work fits perfectly within the scope of JCP as it provides a rigorous
physical chemistry tool that enhances the reliability of structural
parameters. I recommend its publication as it is.}

\paragraph{Response.}
We thank the reviewer for this positive assessment. We are pleased that the
reproducibility, treatment of incomplete data, validation strategy, and clarity
of the presentation were found to be strengths of the work. We have retained
the methodological scope and the benchmark set while making the specific
clarifications and presentation changes requested below.

\subsection*{Comment 1: Computational cost}

\paragraph{Reviewer comment.}
\emph{Computational cost: While the ``testosterone scale check'' is impressive,
a brief mention of the typical computational overhead added by the MORPHEUS
framework compared to traditional manual refinement would be informative for
potential users.}

\paragraph{Response.}
We agree that this information is useful for potential users. We now state
that topology perception, coordinate generation, projection, rank auditing,
and the subsequent refinement are typically completed in seconds on a
laptop-class CPU for the systems considered here, including the 49-atom
testosterone case. We also explain that the sparse Wilson $\mathbf{B}$ matrix
reduces memory use and the cost of the associated linear-algebra operations,
and that no additional quantum-chemical calculation is required by the
model-generation step. The statement is deliberately expressed as a typical
range of practical behavior rather than as a hardware-independent benchmark.

\paragraph{Change in the manuscript.}
A new paragraph was added to the workflow discussion on page~11 of the
revised manuscript.

\subsection*{Comment 2: Figure 1 label}

\paragraph{Reviewer comment.}
\emph{In Figure 1, the label ``algorithmic model generation'' overlaps with an
arrow. I suggest repositioning the text to improve legibility.}

\paragraph{Response.}
We agree and have redesigned the label placement. The label is now set on two
lines and positioned adjacent to the vertical arrow, without covering the
arrow or any other graphical element. This is a graphical correction and is
therefore not marked in red in the marked manuscript.

\paragraph{Change in the manuscript.}
The revised Figure~1 appears on page~3.

\subsection*{Comment 3: Figure 5 size}

\paragraph{Reviewer comment.}
\emph{Figure 5 is currently too small to be legible. I recommend either
increasing the font size or expanding the figure to a double-column format.}

\paragraph{Response.}
We agree. Figure~5 has been expanded to a double-column figure and its labels
and annotations have been enlarged to improve readability at normal manuscript
size. This is a graphical correction and is therefore not marked in red in the
marked manuscript.

\paragraph{Change in the manuscript.}
The revised Figure~5 appears on page~13.

We hope that these revisions address all of the reviewer's suggestions and
improve the clarity and usability of the manuscript.

Sincerely,

Vincenzo Barone

\end{document}
